% Part: first-order-logic
% Chapter: sequent-calculus
% Section: provability-quantifiers

\documentclass[../../../include/open-logic-section]{subfiles}

\begin{document}

\olfileid{fol}{seq}{qpr}
\olsection{\usetoken{S}{derivability} lan Kuantor}

\begin{explain}
  Teorema kelengkapan uga mbutuhake supaya aturan kalkulus sekuen
  ngasilake kasunyatan ngenani~$\Proves$ kang ditetepake ing bagian
  iki.
\end{explain}

\begin{thm}
\ollabel{thm:strong-generalization} Yen $c$ iku konstanta kang ora
muncul ing $\Gamma$ utawa $!A(x)$ lan $\Gamma \Proves !A(c)$, mula
$\Gamma \Proves \lforall[x][!A(x)]$.
\end{thm}

\begin{proof}
Ayo $\pi_0$ dadi $\Log{LK}$-!!{derivation} saka $\Gamma_0 \Sequent !A(c)$
kanggo sawijining $\Gamma_0 \subseteq \Gamma$ kang winates. Kanthi
nambahake inferensi $\RightR{\lforall}$, kita entuk !!a{derivation}
saka $\Gamma_0 \Sequent \lforall[x][!A(x)]$, amarga $c$ ora muncul ing
$\Gamma$ utawa $!A(x)$, mula syarat eigenvariabel wis dicukupi.
\end{proof}

\begin{prop}
\ollabel{prop:provability-quantifiers}
\begin{tagenumerate}{prvEx,prvAll}
\tagitem{prvEx}{$!A(t) \Proves \lexists[x][!A(x)]$.}{}
\tagitem{prvAll}{$\lforall[x][!A(x)] \Proves !A(t)$.}{}
\end{tagenumerate}
\end{prop}

\begin{proof}
\begin{tagenumerate}{prvEx,prvAll}
\tagitem{prvEx}{Sekuen $!A(t) \Sequent \lexists[x][!A(x)]$ bisa
  !!{derivable}:
  \begin{prooftree}
    \Axiom$!A(t) \fCenter !A(t)$
    \RightLabel{\RightR{\lexists}}
    \UnaryInf$!A(t) \fCenter \lexists[x][!A(x)]$
    \end{prooftree}}{}
\tagitem{prvAll}{Sekuen $\lforall[x][!A(x)] \Sequent !A(t)$ bisa
  !!{derivable}:
  \begin{prooftree}
    \Axiom$!A(t) \fCenter !A(t)$
    \RightLabel{\LeftR{\lforall}}
    \UnaryInf$\lforall[x][!A(x)] \fCenter !A(t)$
    \end{prooftree}}{}
\end{tagenumerate}
\end{proof}

\end{document}
